\section{Hiperparametreler}
\label{app:hyper}

Tablo~\ref{tab:hyper}, gömülü hedefe aktarılan modeli üretmek için
kullanılan hiperparametrelerin tamamını listeler. Tüm değerler beş
eğitim tohumu boyunca değişmeden tutulmuştur.

\begin{table}[h]
\centering
\caption{Gömülü hedefe aktarılan \fastgrnn{} modeli için eğitim ve
aktarım hiperparametreleri.}
\label{tab:hyper}
\begin{tabular}{lr}
\toprule
Parametre & Değer \\
\midrule
Gizli boyut $H$                  & 16 \\
Giriş boyutluluğu $d$            & 3 (üç eksenli ivme) \\
Pencere uzunluğu $T$             & 128 örnek (2.56~s) \\
Örnekleme hızı                   & 50~Hz \\
Çıkış sınıfları                  & 6 \\
\midrule
Yinelemeli rank $r_u$            & 8 \\
Giriş rankı $r_w$                & 2 \\
Hedef seyreklik $s$              & 0.5 (283 sıfır olmayan / 566) \\
Seyreklik rampa epoch'u          & 50 (kübik çizelge) \\
Sabit-maskeli ince ayar epoch'u  & 50 \\
\midrule
İyileştirici                     & Adam \\
Öğrenme oranı $\eta$             & $10^{-3}$ \\
Yığın boyutu                     & 64 \\
Toplam epoch                     & 100 \\
Rastgele tohumlar                & $\{0, 1, 2, 3, 4\}$ \\
\midrule
Aktivasyon fonksiyonu            & sigmoid, tanh \\
LUT boyutu                       & her biri 256 giriş \\
LUT giriş alanı                  & $[-8, +8]$ \\
Ağırlık nicelendirme             & tensör-başına Q15 \\
Aktivasyon nicelendirme          & tensör-başına Q15 + kalibrasyon \\
Kalibrasyon mini-yığınları       & 5 \\
Kalibrasyon güvenlik payı        & 10\% \\
\bottomrule
\end{tabular}
\end{table}

\section{Q15 Nicelendirme Uygulaması}
\label{app:q15}

Bölüm~\ref{sec:method:q15}'de kullanılan tensör-başına Q15
nicelendirme adımı, her tensör için tek satırlık bir işlemdir.
Python'da:

\begin{verbatim}
def to_q15(W: np.ndarray, scale: float) -> np.ndarray:
    return (W / scale).round() \
                      .clip(-32768, 32767) \
                      .astype(np.int16)
\end{verbatim}

Tensör-başına ölçek
\[
s_\ell = \max_{ij}\bigl|W^{(\ell)}_{ij}\bigr| \,\big/\, 32767,
\]
olarak seçilir; bu, her nicelendirilmiş girdiyi işaretli 16-bit
aralık içinde tutan en büyük ölçektir. Aynı ifade hem yinelemeli
ağırlık çarpanları $\mathbf{W}_1, \mathbf{W}_2,
\mathbf{U}_1, \mathbf{U}_2$ hem de sınıflandırıcı başlığı için
kullanılır. Çalışma zamanında çözülmüş değer şu şekilde hesaplanır:

\begin{verbatim}
float w = (float)W_q15[i] * scale;
\end{verbatim}

Her çözme işleminin çalışma zamanı maliyeti bir int-to-float
dönüşümü ve bir çarpmadır. Bunlar AVR üzerinde ucuzdur; LUT
aktivasyonları araya girdiğinde çarpıcısız \msp{} üzerinde de
yönetilebilir düzeydedir.

\section{LUT Üretim Algoritması}
\label{app:lut}

Bölüm~\ref{sec:method:lut}'deki sigmoid ve tanh arama tabloları
Python'da çevrim dışı olarak önceden hesaplanır ve C başlık dosyası
olarak yayılır:

\begin{verbatim}
LUT_SIZE     = 256
INPUT_MIN    = -8.0
INPUT_MAX    =  8.0
BUCKET_WIDTH = (INPUT_MAX - INPUT_MIN) / LUT_SIZE

sigmoid_lut = [
    1.0 / (1.0 + math.exp(-(INPUT_MIN
            + (i + 0.5) * BUCKET_WIDTH)))
    for i in range(LUT_SIZE)
]
tanh_lut = [
    math.tanh(INPUT_MIN + (i + 0.5) * BUCKET_WIDTH)
    for i in range(LUT_SIZE)
]
\end{verbatim}

Her giriş, kendi kovasının \emph{merkezindeki} aktivasyon değerini
saklar ($(i + 0.5)$ ofseti). Bu, kovacık içinde düzgün dağılmış bir
giriş için maksimum olabilirlik tahminidir ve yarım kova yanlılığını
önler. Çalışma zamanı araması şöyledir:

\begin{verbatim}
static inline float lut_eval(const float *lut, float x) {
    if (x <= LUT_INPUT_MIN) return lut[0];
    if (x >= LUT_INPUT_MAX) return lut[LUT_SIZE - 1];
    int idx = (int)((x - LUT_INPUT_MIN) * LUT_INPUT_SCALE);
    return lut[idx];
}
\end{verbatim}

\texttt{LUT\_INPUT\_SCALE = 1 / BUCKET\_WIDTH} olarak tanımlanır.
İki tablo birlikte 2~KB Flash kaplar; birlikte çalışma zamanındaki
tüm \texttt{expf} ve \texttt{tanhf} çağrılarını ortadan kaldırırlar.

\section{MSP430 \texorpdfstring{I\textsuperscript{2}C}{I2C} İlk Çalıştırma}
\label{app:hw}

\msp{} üzerinde canlı MPU-6050 okuması, standart pin atamasıyla
(P1.6~$=$~SCL, P1.7~$=$~SDA) I\textsuperscript{2}C ana kipindeki
USCI\_B0 çevre birimini kullanır. Gelecekteki okurlar için iki
pratik hususu kaydetmeye değer:

\begin{itemize}
\item \textbf{J5 jumper'ının çıkarılması.} MSP-EXP430G2 LaunchPad,
kart üzerindeki yeşil LED'i P1.6 ile çoklar. Herhangi bir
I\textsuperscript{2}C iletişiminden önce J5 jumper'ı fiziksel
olarak çıkarılmalıdır; aksi hâlde SCL LED tarafından yüklenir ve
veri yolu I\textsuperscript{2}C yüksek seviye eşiğinin altında kalır.

\item \textbf{Bus-busy kurtarma.} Aktarım ortasında güç kesilirse
USCI\_B0 çevre birimi \texttt{UCSCLLOW + UCBBUSY} durumuna
kilitlenebilir. Geleneksel 9-saat kurtarma dizisi (SDA serbestken
SCL'yi dokuz kez değiştirip ardından STOP üretmek) veri yolunu
geri yükler. Açılışta GPIO düzeyinde yapılan bir sağlık denetimi
(\texttt{P1IN \& (BIT6|BIT7)} değerinin \texttt{0xC0} olması
beklenir), I\textsuperscript{2}C durum makinesi devreye alınmadan
önce bu koşulu saptar.
\end{itemize}

Bu iki madde yaklaşık bir günlük ilk çalıştırma süresi tüketmiştir
ve standart MSP430 uygulama notlarında~\cite{ti_msp430}
belgelenmemiştir. Bir sonraki okura aynı günü kazandırma olasılığı
için buraya eklenmiştir.
